% mission chez Orange
%%%%%%%%%%%%%%%%%%%%%%%%%%%%%%%%%%%%%%%%%%%%%%%%%%%%%%%%%%%%%%%%%%%%%%%%%%%%%%%%

\pagestyle{highlightmain}
\highlightbar{
    \section{Extia}
    \textit{09/2018 - 12/2020}\\
    \\Extia est une société française de conseil et d'ingénierie informatique proposant des prestations de développement logiciel et des services numériques dans de nombreux secteurs.
    \\Elle intervient en mission chez ses clients pour les accompagner dans la conception, le développement et la maintenance d'applications métier.
    \vspace{0.5em}\\
    \location{Aix-en-Provence, France}
    \homepage{extia-group.com}{https://www.extia-group.com/}
    \linkedin{Extia}{https://www.linkedin.com/company/extia/}
    \vspace{0.5em}\\
    Référent :\\
    \contact{Cédric VINCENT, Directeur d'agence}{cvincent@extia.fr}

    \section{Orange (\textit{Mission})}
    \textit{09/2020 - 12/2020}\\
    \\Orange est un groupe français de télécommunications proposant des services mobiles, fixes, d'accès à Internet haut débit et des services numériques.
    \\Il accompagne des millions de clients en France et dans le monde avec des infrastructures réseau, des solutions numériques et des services informatiques pour les particuliers et les entreprises.
    \vspace{0.5em}\\
    \location{Marseille, France}
    \homepage{orange.fr}{https://www.orange.fr/}
}
\mainbar{
    \section[\faGears]{Extia - Consultant full-stack}
    \textit{09/2020 - 12/2020 (mission chez Orange)}

    Contributions techniques aux applications de tarification et de gestion de catalogue d'Orange, utilisées pour gérer les offres télécoms, les règles tarifaires et les traitements batch.

    \subsection{Principales responsabilités}
    \begin{itemize}
        \item Contribution à la maintenance et à l'évolution des applications internes de tarification et de gestion de catalogue d'Orange.
        \\\item Développement de fonctionnalités back-end et front-end avec \textcolor{accent}{Spring Boot} et \textcolor{accent}{Angular} au sein d'une équipe produit full-stack.
        \\\item Conception et développement d'un tableau de bord de supervision des traitements \textcolor{accent}{Spring Batch}, avec \textcolor{accent}{Spring Actuator} et Material, pour un suivi opérationnel en temps réel.
        \\\item Développement d'API d'orchestration et de fonctions de supervision permettant aux équipes d'exploitation de suivre et de piloter les traitements batch de l'écosystème produit.
        \\\item Réalisation d'un prototype full-stack associant tableaux de bord Angular, mécanismes de rafraîchissement et endpoints de suivi des batchs pour valider l'approche technique de la plateforme de supervision.
        \\\item Travail en autonomie sur la supervision transverse, après avoir rapidement gagné la confiance du Tech Lead et de l'équipe projet.
        \\\item Collaboration avec les équipes de qualification pour valider les fonctionnalités et préparer les livraisons.
    \end{itemize}

    \subsection{Technologies \& outils}
    \tag{Java} \tag{Spring Boot}[web, security, actuator, batch, data]\\
    \tag{Angular} \tag{Material}\\
    \tag{API REST} \tag{Orchestration batch} \tag{Supervision}\\
    \tag{IntelliJ} \tag{Docker} \tag{GitLab} \tag{Oracle}

    \subsection{Langues utilisées}
    \tag{Français}[quotidien, langue maternelle]
}
\makebody
\clearpage

% mission chez AVIS2SANTÉ
%%%%%%%%%%%%%%%%%%%%%%%%%%%%%%%%%%%%%%%%%%%%%%%%%%%%%%%%%%%%%%%%%%%%%%%%%%%%%%%%

\pagestyle{highlightmain}
\highlightbar{
    \section{Extia}
    \textit{09/2018 - 12/2020}\\
    \\Extia est une société française de conseil et d'ingénierie informatique proposant des prestations de développement logiciel et des services numériques dans de nombreux secteurs.
    \\Elle intervient en mission chez ses clients pour les accompagner dans la conception, le développement et la maintenance d'applications métier.
    \vspace{0.5em}\\
    \location{Aix-en-Provence, France}
    \homepage{extia-group.com}{https://www.extia-group.com/}
    \linkedin{Extia}{https://www.linkedin.com/company/extia/}
    \vspace{0.5em}\\
    Référent :\\
    \contact{Cédric VINCENT, Directeur d'agence}{cvincent@extia.fr}

    \section{AVIS2SANTÉ (\textit{Mission})}
    \textit{04/2020 - 08/2020}\\
    \\AVIS2SANTÉ est une PME française de la santé numérique développant des logiciels de télémédecine et de santé connectée pour les professionnels de santé et les patients.
    \\Elle conçoit des outils web sécurisés de gestion des rendez-vous, de consultation médicale et de partage d'informations entre praticiens et patients pour accompagner les services de santé numérique.
    \vspace{0.5em}\\
    \location{Six-Fours-les-Plages, France}
    \homepage{avis2sante.fr}{https://avis2sante.fr/}
}
\mainbar{
    \section[\faGears]{Extia - Consultant full-stack}
    \textit{04/2020 - 08/2020 (mission chez AVIS2SANTÉ)}

    Contributions techniques à plusieurs projets web et mobiles chez AVIS2SANTÉ (A2S), PME de la santé : évolution de plateformes web, intégration d'objets connectés et amélioration des pratiques de développement pendant la période COVID.

    \subsection{Principales responsabilités}
    \begin{itemize}
        \item Développement de modules Joomla et d'intégrations pour des plateformes web de santé, dont la prise de rendez-vous et des modules de paiement \textcolor{accent}{Stripe}.
        \\\item Mise en œuvre d'intégrations front-end JavaScript réutilisables dans Joomla, avec des bibliothèques légères et des architectures de modules spécifiques.
        \\\item Contribution à un prototype Android de synchronisation de données de santé connectée avec \textcolor{accent}{Kotlin} et Android Studio.
        \\\item Intégration des API \textcolor{accent}{Withings} et \textcolor{accent}{Google Fit} pour synchroniser les données de suivi de santé entre les objets connectés et les plateformes A2S.
        \\\item Participation à la conception des flux de synchronisation mobile, du stockage local et des échanges avec les services back-end Joomla.
        \\\item Accompagnement technique d'un développeur Android/Kotlin en alternance et participation aux échanges sur l'architecture de l'application mobile.
        \\\item Collaboration avec le CTO pour améliorer les processus de développement, notamment la gestion de versions avec \textcolor{accent}{GitLab} et les pratiques de développement partagées.
        \\\item Travail en grande autonomie dans une petite équipe pluridisciplinaire, avec présentation régulière de l'avancement et démonstration de prototypes aux interlocuteurs techniques et métier.
    \end{itemize}

    \subsection{Technologies \& outils}
    \tag{PHP} \tag{Joomla} \tag{JavaScript} \tag{Bibliothèques de type jQuery}\\
    \tag{Kotlin} \tag{Android Studio} \tag{Google Fit} \tag{API Withings}\\
    \tag{API Stripe} \tag{GitLab} \tag{Docker}

    \subsection{Langues utilisées}
    \tag{Français}[quotidien, langue maternelle]
}
\makebody
\clearpage

% mission chez Monext
%%%%%%%%%%%%%%%%%%%%%%%%%%%%%%%%%%%%%%%%%%%%%%%%%%%%%%%%%%%%%%%%%%%%%%%%%%%%%%%%

\pagestyle{highlightmain}
\highlightbar{
    \section{Extia}
    \textit{09/2018 - 12/2020}\\
    \\Extia est une société française de conseil et d'ingénierie informatique proposant des prestations de développement logiciel et des services numériques dans de nombreux secteurs.
    \\Elle intervient en mission chez ses clients pour les accompagner dans la conception, le développement et la maintenance d'applications métier.
    \vspace{0.5em}\\
    \location{Aix-en-Provence, France}
    \homepage{extia-group.com}{https://www.extia-group.com/}
    \linkedin{Extia}{https://www.linkedin.com/company/extia/}
    \vspace{0.5em}\\
    Référent :\\
    \contact{Cédric VINCENT, Directeur d'agence}{cvincent@extia.fr}

    \section{Monext (\textit{Mission})}
    \textit{12/2018 - 03/2020}\\
    \\Monext est un prestataire français de services de paiement proposant des solutions de paiement électronique, de gestion de la fraude et de commerce numérique.
    \\L'entreprise exploite des plateformes sécurisées pour les transactions en ligne et en magasin des commerçants, banques et établissements financiers.
    \vspace{0.5em}\\
    \location{Aix-en-Provence, France}
    \homepage{monext.com}{https://www.monext.com/}
}
\mainbar{
    \section[\faGears]{Extia - Consultant full-stack}
    \textit{12/2018 - 03/2020 (mission chez Monext)}

    Contributions techniques à la plateforme de paiement PayAvenue chez Monext : traitements d'acquisition à fort volume et intégration de protocoles bancaires pour de grandes enseignes.

    \subsection{Principales responsabilités}
    \begin{itemize}
        \item Développement et maintenance de traitements \textcolor{accent}{Spring Batch} pour les services d'acquisition de paiements traitant les transactions de grandes enseignes à fort volume.
        \\\item Contribution aux flux d'acquisition Carrefour avec les protocoles \textcolor{accent}{NEXO}, pour accompagner la transition depuis les formats bancaires historiques \textcolor{accent}{CB2A} et \textcolor{accent}{CB2C}.
        \\\item Participation aux traitements multi-acquisitions permettant plusieurs échanges quotidiens à la place d'un unique cycle de traitement en fin de journée.
        \\\item Implémentation de la conversion de protocoles et des intégrations bancaires entre les systèmes de paiement des enseignes et plusieurs banques acquéreuses.
        \\\item Prise en compte des contraintes critiques du paiement : plages d'acquisition, agrégation des transactions, sécurisation des échanges et volumes élevés.
        \\\item Contribution à la maintenance, aux analyses techniques et à l'évolution des services de paiement en environnement \textcolor{accent}{PCI DSS}.
        \\\item Participation à l'intégration, au support de qualification et à la préparation des livraisons des applications de paiement.
    \end{itemize}

    \subsection{Technologies \& outils}
    \tag{Java} \tag{Spring Boot}[batch, data] \tag{Maven}\\
    \tag{Protocoles de paiement}[CB2A, NEXO]\\
    \tag{Sécurité}[PCI DSS, hachage, RSA]\\
    \tag{Jenkins} \tag{Docker} \tag{SVN} \tag{Tortoise}\\
    \tag{Oracle} \tag{Eclipse}

    \subsection{Langues utilisées}
    \tag{Français}[quotidien, langue maternelle]
}
\makebody
\clearpage

% mission chez Docapost
%%%%%%%%%%%%%%%%%%%%%%%%%%%%%%%%%%%%%%%%%%%%%%%%%%%%%%%%%%%%%%%%%%%%%%%%%%%%%%%%

\pagestyle{highlightmain}
\highlightbar{
    \section{Extia}
    \textit{09/2018 - 12/2018}\\
    \\Extia est une société française de conseil et d'ingénierie informatique proposant des prestations de développement logiciel et des services numériques dans de nombreux secteurs.
    \\Elle intervient en mission chez ses clients pour les accompagner dans la conception, le développement et la maintenance d'applications métier.
    \vspace{0.5em}\\
    \location{Aix-en-Provence, France}
    \homepage{extia-group.com}{https://www.extia-group.com/}
    \linkedin{Extia}{https://www.linkedin.com/company/extia/}
    \vspace{0.5em}\\
    Référent :\\
    \contact{Cédric VINCENT, Directeur d'agence}{cvincent@extia.fr}

    \section{Docapost (\textit{Mission})}
    \textit{09/2018 - 12/2018}\\
    \\Docaposte est la filiale numérique du Groupe La Poste, proposant des solutions sécurisées de confiance et de transformation numériques aux entreprises et aux institutions publiques.
    \\Elle fournit des services de signature électronique, d'archivage numérique, d'échange sécurisé de données et d'intégration de plateformes pour les processus réglementés et les transformations numériques à grande échelle, notamment dans les services publics, la santé et la finance.
    \\\textcolor{accent!90}{\textit{Le Groupe La Poste est un acteur français majeur des services postaux et logistiques, proposant des services de courrier, de colis, bancaires et numériques en France et à l'international.
    \\Il assure des missions de service public en France auprès de millions de clients et accompagne les particuliers, les entreprises et les institutions dans de nombreux secteurs.}}
    \vspace{0.5em}\\
    \location{Aix-en-Provence, France}
    \homepage{docaposte.com}{https://www.docaposte.com/}
}
\mainbar{
    \section[\faGears]{Extia - Consultant full-stack}
    \textit{09/2018 - 12/2018 (mission chez Docapost)}

    Contributions techniques à la plateforme d'immatriculation de Docapost pour La Poste : dématérialisation des démarches et gestion documentaire interne après l'ouverture du marché français de l'immatriculation.

    \subsection{Principales responsabilités}
    \begin{itemize}
        \item Contribution aux applications web internes de gestion des immatriculations et des documents au sein des services BPO de Docapost.
        \\\item Création d'une application back-office avec \textcolor{accent}{Spring MVC} et Bootstrap pour aider le support de niveau 3 à suivre et à analyser les dossiers d'immatriculation.
        \\\item Intégration d'\textcolor{accent}{API REST} avec les services back-end de référence centralisés, utilisés par les applications web et mobiles du projet.
        \\\item Développement de vues front-end et de contrôleurs back-end avec Spring MVC, \textcolor{accent}{JSP} et Bootstrap dans un environnement Java d'entreprise.
        \\\item Collaboration avec les développeurs back-end seniors pour faire évoluer les API de référence et ajouter les endpoints nécessaires aux processus opérationnels.
        \\\item Prise en charge en grande autonomie de besoins urgents d'outillage opérationnel peu après l'arrivée dans l'équipe projet.
        \\\item Participation au support de qualification, à l'intégration et à la préparation des déploiements sur des environnements \textcolor{accent}{Tomcat}.
    \end{itemize}

    \subsection{Technologies \& outils}
    \tag{Java} \tag{Spring Framework}[web, security, data]\\
    \tag{API REST} \tag{Swagger} \tag{Bootstrap} \tag{JSP}\\
    \tag{Tomcat} \tag{Git} \tag{Jenkins} \tag{Eclipse}

    \subsection{Langues utilisées}
    \tag{Français}[quotidien, langue maternelle]
}
\makebody
\clearpage